% Figure : carte de préconisation par inférence à rebours sur le réseau appris, en nuage de points.
% Probabilité a posteriori de chaque profil sachant IPC = faible, pour les sept contextes (nature, position de la cible),
% un plan par durée du choc. Vert : politique la plus probable du contexte ; rouge : autres politiques.
% Valeurs identiques à l'ancienne carte (inférence exacte sur le réseau appris, fichier .bif).
\begin{tikzpicture}
\pgfplotsset{carte/.style={width=5.6cm, height=5.6cm, xmin=0.5, xmax=4.5, ymin=0, ymax=0.85,
  xtick={1,2,3,4}, xticklabels={Aucune,Tardif,Prudent,Réactif},
  ytick={0,0.2,0.4,0.6,0.8}, yticklabel style={/pgf/number format/.cd,fixed,fixed zerofill,precision=1,use comma},
  tick label style={font=\scriptsize}, label style={font=\scriptsize}, title style={font=\small\bfseries},
  ymajorgrids, grid style={gray!20}, axis lines*=left, clip=false}}
\begin{axis}[carte, title={Durée courte}, at={(0.0cm,0)}, ylabel={Probabilité a posteriori du profil}]
  \addplot[only marks, mark=*, mark size=1.8pt, appreac, fill opacity=0.75, draw opacity=0.9] coordinates {(0.73,0.16) (1.73,0.15) (2.73,0.12) (0.82,0.15) (1.82,0.14) (2.82,0.13) (0.91,0.13) (1.91,0.18) (2.91,0.17) (1.00,0.14) (2.00,0.18) (3.00,0.17) (1.09,0.14) (2.09,0.17) (3.09,0.17) (1.18,0.15) (2.18,0.15) (3.18,0.15) (1.27,0.15) (2.27,0.15) (3.27,0.15)};
  \addplot[only marks, mark=*, mark size=1.8pt, apptard, fill opacity=0.85] coordinates {(3.73,0.58) (3.82,0.57) (3.91,0.52) (4.00,0.51) (4.09,0.52) (4.18,0.55) (4.27,0.55)};
\end{axis}
\begin{axis}[carte, title={Durée moyenne}, at={(5.15cm,0)}, yticklabels={}]
  \addplot[only marks, mark=*, mark size=1.8pt, appreac, fill opacity=0.75, draw opacity=0.9] coordinates {(0.73,0.07) (1.73,0.07) (2.73,0.10) (0.82,0.06) (1.82,0.07) (2.82,0.12) (0.91,0.06) (1.91,0.10) (2.91,0.14) (1.00,0.06) (2.00,0.10) (3.00,0.14) (1.09,0.06) (2.09,0.10) (3.09,0.14) (1.18,0.07) (2.18,0.08) (3.18,0.12) (1.27,0.07) (2.27,0.08) (3.27,0.12)};
  \addplot[only marks, mark=*, mark size=1.8pt, apptard, fill opacity=0.85] coordinates {(3.73,0.76) (3.82,0.75) (3.91,0.69) (4.00,0.70) (4.09,0.70) (4.18,0.73) (4.27,0.73)};
\end{axis}
\begin{axis}[carte, title={Durée longue}, at={(10.3cm,0)}, yticklabels={}]
  \addplot[only marks, mark=*, mark size=1.8pt, appreac, fill opacity=0.75, draw opacity=0.9] coordinates {(0.73,0.06) (1.73,0.07) (2.73,0.08) (0.82,0.06) (1.82,0.07) (2.82,0.10) (0.91,0.05) (1.91,0.10) (2.91,0.13) (1.00,0.05) (2.00,0.09) (3.00,0.13) (1.09,0.05) (2.09,0.09) (3.09,0.13) (1.18,0.06) (2.18,0.08) (3.18,0.10) (1.27,0.06) (2.27,0.08) (3.27,0.10)};
  \addplot[only marks, mark=*, mark size=1.8pt, apptard, fill opacity=0.85] coordinates {(3.73,0.79) (3.82,0.77) (3.91,0.72) (4.00,0.72) (4.09,0.73) (4.18,0.76) (4.27,0.76)};
\end{axis}
\matrix[anchor=north, font=\scriptsize, column sep=4pt] at (7.6cm,-0.5cm) {
  \fill[apptard] (0,0) circle (2pt); & \node[anchor=west]{Politique la plus probable du contexte}; &
  \fill[appreac] (0,0) circle (2pt); & \node[anchor=west]{Autres politiques}; \\
};
\end{tikzpicture}
